\section{Loop structure and harmonic machinery}
\label{sec:harmonic}
This section turns direction-space kernels into $C^{(rs)}_\ell$.  It acts on the maps of Secs.~\ref{sec:linear}--\ref{sec:map3}
through the normalisation algebra of Sec.~\ref{sec:normalisation}; all rules are in the single convention
\eqref{eq:bgray}--\eqref{eq:rayleigh} (phase $e^{+i\vk\cdot\nhat\chi}$, factor $i^\ell$).

\subsection{Loop structure}
\label{sec:loops}
Write each order as a direction-space kernel acting on the Gaussian field,
\begin{equation}
\Theta^{(r)}(\nhat)=\int\prod_{i=1}^r\frac{\dd^3k_i}{(2\pi)^3}\,\zeta_{\vk_1}\cdots\zeta_{\vk_r}\,K_r(\vk_1,\ldots,\vk_r;\nhat),
\qquad K_{r,\ell m}(\vk_1,\ldots)\equiv\int\!\dd\Omega_{\nhat}Y^*_{\ell m}(\nhat)K_r ,
\label{eq:Kr}
\end{equation}
with $K_r$ symmetric in its momenta (symmetrise if necessary).  Wick's theorem for $\ell\ge1$ then gives
\begin{align}
C^{(11)}_\ell&=\frac1{2\ell+1}\sum_m\int\!\frac{\dd^3k}{(2\pi)^3}P(k)\,\big|K_{1,\ell m}(\vk)\big|^2,\nonumber\\
C^{(22)}_\ell&=\frac{2}{2\ell+1}\sum_m\int\!\frac{\dd^3k_1\,\dd^3k_2}{(2\pi)^6}P(k_1)P(k_2)\,\big|K_{2,\ell m}(\vk_1,\vk_2)\big|^2,\nonumber\\
C^{(13)}_\ell&=\frac{3}{2\ell+1}\sum_m\int\!\frac{\dd^3k}{(2\pi)^3}P(k)\,K_{1,\ell m}(\vk)\Big[\int\!\frac{\dd^3q}{(2\pi)^3}P(q)\,K_{3,\ell m}(\vk,\vq,-\vq)\Big]^* .
\label{eq:loops}
\end{align}
The factor 2 in $C^{(22)}$ counts the two connected pairings of a symmetric $K_2$; the third pairing
$\langle\zeta_{\vk_1}\zeta_{\vk_2}\rangle$ inside $\Theta^{(2)}$ has $\vk_1+\vk_2=0$ and feeds only $\ell=0$.
The factor 3 in $C^{(13)}$ counts the three choices of which $\zeta$ inside $\Theta^{(3)}$ pairs with
$\Theta^{(1)}$; the remaining pair is contracted at $\vq_j=-\vq_k$, so the ``tadpole'' configuration
$K_3(\vk,\vq,-\vq)$ \emph{is} the entire $(13)$ term and must be kept (this is where the mean-field terms of
Sec.~\ref{sec:third} enter).  Both prefactors have been re-derived independently and confirmed on the lensing
control case below (App.~\ref{app:ledger}, H-6).  Equation~\eqref{eq:loops} uses $\langle a^{(1)}a^{(3)*}\rangle$;
only $\Re C^{(13)}_\ell$ enters \eqref{eq:cl_target}.  With the contracted fields of \eqref{eq:contracted}, $3\int P(q)K_3(\vk,\vq,-\vq)$ is the
$q$-contracted map $\Theta^{(3)}_c(\vk;\nhat)$ built from $X_{3c}$ and the contraction rules
\eqref{eq:contraction_rules}.

\subsection{Direction convention}
\label{sec:dirconv}
A field on the background ray \eqref{eq:bgray} (observer at $\bar\vx_o=0$) reads
$f(\eta,\nhat\chi)=\int\frac{\dd^3k}{(2\pi)^3}f_{\vk}\,e^{+i\vk\cdot\nhat\chi}$, and all rules below follow from the Rayleigh expansion
\eqref{eq:rayleigh}.  Because $Y_{\ell m}(-\nhat)=(-1)^\ell Y_{\ell m}(\nhat)$, expressing the same maps through the propagation
direction $-\nhat$ (phase $e^{-i\vk\cdot(-\nhat)\chi}$, factors $(-i)^\ell$) multiplies every $a_{\ell m}$ by $(-1)^\ell$ at every order;
the transfer functions $\Delta_\ell$ (the factor $i^\ell$ versus $(-i)^\ell$ absorbs the sign) and all $C^{(rs)}_\ell$ are unchanged.  One
convention must be used for $\Theta^{(1)},\Theta^{(2)},\Theta^{(3)}$ together; we use $\nhat$ throughout.  Gaunt integrals are
\begin{equation}
\mathcal G^{\ell\,\ell_1\ell_2}_{m\,m_1m_2}\equiv\int\!\dd\Omega\,Y^*_{\ell m}Y_{\ell_1m_1}Y_{\ell_2m_2}
=(-1)^m\sqrt{\frac{(2\ell+1)(2\ell_1+1)(2\ell_2+1)}{4\pi}}\threej{\ell}{\ell_1}{\ell_2}{0}{0}{0}\threej{\ell}{\ell_1}{\ell_2}{-m}{m_1}{m_2}.
\end{equation}

\subsection{Type A: linear line-of-sight maps}
\label{sec:typeA}
A map of the form
\begin{equation}
\Theta_A(\nhat)=\int\!\frac{\dd^3k}{(2\pi)^3}\zeta_{\vk}\int_{\eta_*}^{\bar\eta_o}\!\dd\eta\sum_{s\ge0}g_s(k,\eta)\,(i\,\nhat\cdot\khat)^s\,e^{i\vk\cdot\nhat\chi(\eta)}
\label{eq:typeA}
\end{equation}
(endpoint terms are $\delta(\eta-\eta_1)$ pieces of $g_s$; the monomials are written as $(i\nhat\cdot\khat)^s$
because $i\khat=\nabla/k$, so a real map has real $g_s$) has
\begin{equation}
\text{(A-$s$)}\qquad a_{\ell m}=4\pi\,i^\ell\int\!\frac{\dd^3k}{(2\pi)^3}\zeta_{\vk}Y^*_{\ell m}(\khat)\,\Delta_\ell(k),\qquad
\Delta_\ell(k)=\int\!\dd\eta\sum_sg_s(k,\eta)\,\jl^{(s)}\big(k\chi(\eta)\big),
\label{eq:ruleA}
\end{equation}
i.e.\ $(i\nhat\cdot\khat)^se^{i\vk\cdot\nhat\chi}\to\jl^{(s)}(x)$, $x=k\chi$, since $(i\nhat\cdot\khat)e^{ik\chi\nhat\cdot\khat}=k^{-1}\partial_\chi e^{ik\chi\nhat\cdot\khat}$.
Helicity pieces (vectors $\omega_i,v^T_i$ contracted with $\nhat$, tensors $h_{ij}n^in^j$): in the frame
$\khat=\hat{\mathbf z}$ with $e^{(\pm1)}=(\hat{\mathbf x}\pm i\hat{\mathbf y})/\sqrt2$, $e^{(\pm2)}_{ij}=e^{(\pm1)}_ie^{(\pm1)}_j$,
$Q_{\pm1}=\nhat\cdot e^{(\pm1)}$, $Q_{\pm2}=e^{(\pm2)}_{ij}n^in^j$,
\begin{equation}
\begin{split}\text{(A-$\lambda$)}\qquad Q_\lambda e^{i\vk\cdot\nhat\chi}&=\sum_\ell i^\ell\sqrt{4\pi(2\ell+1)}\,\hat\Delta^{(\lambda)}_\ell(x)\,Y_{\ell\lambda}(\nhat),\\
\hat\Delta^{(\pm1)}_\ell&=\pm i\sqrt{\tfrac{\ell(\ell+1)}2}\frac{\jl(x)}x,\quad
\hat\Delta^{(\pm2)}_\ell=-\tfrac12\sqrt{\tfrac{(\ell+2)!}{(\ell-2)!}}\frac{\jl(x)}{x^2},\end{split}
\label{eq:ruleAlambda}
\end{equation}
and in a general frame $a^{(\lambda)}_{\ell m}=4\pi\,i^\ell\int\frac{\dd^3k}{(2\pi)^3}\zeta^{(\lambda)}_{\vk}(-1)^\lambda[{}_{-\lambda}Y_{\ell m}(\khat)]^*\Delta^{(\lambda)}_\ell(k)$,
$\Delta^{(\lambda)}_\ell=\int\dd\eta\,g^{(\lambda)}\hat\Delta^{(\lambda)}_\ell(k\chi)$.  (Relative to the propagation-direction convention only
$\hat\Delta^{(\pm1)}$ changes sign, because $Q_{\pm1}$ is odd in $\nhat$.)  Scalar, vector and tensor pieces
of the same $\vk$ do not cross-correlate.  For two type-A maps
\begin{equation}
\begin{gathered}C^{AA'}_\ell=\frac2\pi\int k^2\dd k\,P(k)\,\Delta^A_\ell(k)\Delta^{A'}_\ell(k)^*,\qquad
\sky A=\int\!\frac{\dd^3k}{(2\pi)^3}\zeta_{\vk}\Delta^A_0(k),\\
\big\langle\sky A\sky{A'}\big\rangle=\frac1{2\pi^2}\int k^2\dd k\,P\,\Delta^A_0\Delta^{A'}_0=\frac{C^{AA'}_0}{4\pi},\end{gathered}
\label{eq:C11}
\end{equation}
and $\langle\sky A\,a^{A'}_{\ell m}\rangle=0$ for $\ell\ge1$ by isotropy.  The first-order map $m_1=L_1$ of
\eqref{eq:L1} is of type A (with $s\le1$ and endpoint terms), so $C^{(11)}_\ell$ is \eqref{eq:C11} with
$A=A'=m_1$ (closed form: Sec.~\ref{sec:linear}, Proposition~\ref{prop:Delta}).

\paragraph{Same-point correlators in closed form.}  The rules below need the zero-lag correlator of two type-A maps,
$\xi^{XY}(0)\equiv\langle X(\nhat)Y(\nhat)\rangle=\sum_\ell\frac{2\ell+1}{4\pi}C^{XY}_\ell$.  The $\ell$ sum need not be performed: writing
$x(k,\nu)\equiv\int\dd\eta\sum_sg^X_s(k,\eta)(i\nu)^se^{ik\nu\chi(\eta)}$ for the plane-wave profile of $X$ ($\nu=\nhat\cdot\khat$), the
Rayleigh expansion gives $x(k,\nu)=\sum_\ell(2\ell+1)i^\ell\Delta^X_\ell(k)P_\ell(\nu)$, and Legendre orthogonality (the addition theorem,
$\sum_\ell(2\ell+1)\jl(z)^2=1$, in its general form) yields
\begin{equation}
\xi^{XY}(0)=\int\!\frac{\dd^3k}{(2\pi)^3}P(k)\;\frac12\int_{-1}^1\!\dd\nu\;x(k,\nu)\,y(k,\nu)^*
=\frac1{2\pi^2}\int k^2\dd k\,P(k)\sum_\ell(2\ell+1)\Delta^X_\ell\Delta^{Y*}_\ell ,
\label{eq:xi0}
\end{equation}
a single $\nu$ integral of explicitly known profiles (for endpoint terms elementary, for line-of-sight terms a double time integral
of $e^{ik\nu(\chi_1-\chi_2)}$, whose $\nu$ integral is $2j_0(k(\chi_1-\chi_2))$ times polynomials).  The angular-gradient version
$\xi^{\nabla X\cdot\nabla Y}(0)=\sum_\ell\frac{\ell(\ell+1)(2\ell+1)}{4\pi}C^{XY}_\ell$ is the same with $x\,y^*$ replaced by
$(1-\nu^2)\,\partial_\nu x\,\partial_\nu y^*$ (since $\nabla_{\nhat}f(\nhat\cdot\khat)=f'(\nu)(\khat-\nu\nhat)$).

\subsection{Type B: products of maps at the same point}
\label{sec:typeB}
For $B=AA'$ and $B'=CD$ with $A,A',C,D$ type-A maps, $b_{\ell m}=\sum\mathcal G^{\ell\ell_1\ell_2}_{mm_1m_2}w(\ell,\ell_1,\ell_2)a_{\ell_1m_1}a'_{\ell_2m_2}$,
where $w=1$ for a plain product and
\begin{equation}
w_\nabla(\ell,\ell_1,\ell_2)=\tfrac12\big[\ell_1(\ell_1+1)+\ell_2(\ell_2+1)-\ell(\ell+1)\big]
\label{eq:wgrad}
\end{equation}
for $\nabla_{\nhat}A\cdot\nabla_{\nhat}A'$ ($w_\nabla$ is symmetric in $\ell_1\leftrightarrow\ell_2$).  Wick's theorem with
$\langle a^X_{\ell m}a^Y_{\ell'm'}\rangle=(-1)^mC^{XY}_\ell\delta_{\ell\ell'}\delta_{m,-m'}$ gives, for $\ell\ge1$,
\begin{equation}
\text{(B1)}\qquad C^{BB'}_\ell=\sum_{\ell_1\ell_2}\frac{(2\ell_1+1)(2\ell_2+1)}{4\pi}\threej{\ell}{\ell_1}{\ell_2}{0}{0}{0}^2
w\,w'\Big[C^{AC}_{\ell_1}C^{A'D}_{\ell_2}+C^{AD}_{\ell_1}C^{A'C}_{\ell_2}\Big],
\label{eq:B1}
\end{equation}
with $\ell+\ell_1+\ell_2$ even; \emph{both} Wick pairings appear and reduce to a single term with a factor 2
only when $A=A'$ (the third pairing $\langle AA'\rangle\langle CD\rangle$ feeds $\ell=0$ only).
Monopole$\times$map terms, with $\sky X\cdot1$ treated as an $\ell=0$ map, give for $\ell\ge1$
\begin{align}
\text{(i)}&\quad \tfrac1{2\ell+1}\textstyle\sum_m\big\langle|\sky X\,y_{\ell m}|^2\big\rangle=\langle\sky X^2\rangle\,C^{YY}_\ell,\nonumber\\
\text{(ii)}&\quad \tfrac1{2\ell+1}\textstyle\sum_m\big\langle(AA')_{\ell m}(\sky X\,y_{\ell m})^*\big\rangle=\langle\sky A\sky X\rangle C^{A'Y}_\ell+\langle\sky{A'}\sky X\rangle C^{AY}_\ell,\nonumber\\
\text{(iii)}&\quad \tfrac1{2\ell+1}\textstyle\sum_m\big\langle(\sky X y_{\ell m})(\sky Z w_{\ell m})^*\big\rangle=\langle\sky X\sky Z\rangle C^{YW}_\ell .
\label{eq:monopole_terms}
\end{align}
Hence, for the type-B part of $m_2$ and with \eqref{eq:theta_orders},
\begin{equation}
\tfrac14C^{(22)}_\ell=\tfrac14\Big[C^{m_2m_2}_\ell-4\,\Re\,\text{(ii)}\big|_{AA'\to m_2,\ X=Y=m_1}+4\langle\sky{m_1}^2\rangle C^{(11)}_\ell\Big].
\label{eq:C22_B}
\end{equation}
\emph{Type D} (nested, unequal-time products such as $\int\dd\eta_1\int_{\eta_1}\dd\eta_2\,W\,(\delta_{ij}-n_in_j)\partial_iX(\eta_1)\partial_jY(\eta_2)$,
the lensing-type terms) reduce to (B1) applied to ``shell maps'' $\tilde X(\nhat)=X(\eta_1,\nhat\chi_1)$ with
unequal-time spectra $C^{xy}_\ell(\eta,\eta')=\frac2\pi\int k^2\dd k\,P\,x(k,\eta)y(k,\eta')^*\jl(k\chi)\jl(k\chi')$, using the exact identity
$(\delta-\nhat\nhat)_{ij}\partial_iX\partial_jY=(\chi_1\chi_2)^{-1}\nabla_{\nhat}\tilde X\cdot\nabla_{\nhat}\tilde Y$ (weight $w_\nabla$) and
$\nhat\cdot\partial X=\partial_\chi\tilde X$ for radial contractions ($\jl\to k\jl'$, weight 1).

\subsection{Type C: convolution kernels}
\label{sec:typeC}
A genuinely second-order map
\begin{equation}
\Theta_C(\nhat)=\int\!\frac{\dd^3k_1\dd^3k_2}{(2\pi)^6}\zeta_{\vk_1}\zeta_{\vk_2}\int\!\dd\eta\,K(k_1,k_2,K,\eta)\,\mathcal Q(\nhat;\khat_1,\khat_2,\Khat)\,e^{i\vK\cdot\nhat\chi},
\qquad \vK=\vk_1+\vk_2,
\label{eq:typeC}
\end{equation}
with $\mathcal Q=\sum q_{abc}(k_1,k_2,\mu)(\nhat\cdot\khat_1)^a(\nhat\cdot\khat_2)^b(\nhat\cdot\Khat)^c$ a polynomial of degree $d$
(transverse and TT projectors contracted with $\nhat$ are of this form), has
\begin{equation}
K_{2,\ell m}(\vk_1,\vk_2)=\int\!\dd\eta\,K\sum_{L=|\ell-d|}^{\ell+d}\sum_M4\pi\,i^L\,j_L(K\chi)\,Y^*_{LM}(\Khat)\,\langle\ell m|\mathcal Q|LM\rangle,
\qquad \langle\ell m|\mathcal Q|LM\rangle=\int\!\dd\Omega\,Y^*_{\ell m}\mathcal Q\,Y_{LM},
\label{eq:K2lm}
\end{equation}
a finite $3j$ sum generated by repeated use of $(\nhat\cdot\hat{\mathbf v})Y_{\ell m}=\frac{4\pi}3\sum_{m''}Y^*_{1m''}(\hat{\mathbf v})\sum_{\ell'=\ell\pm1,m'}\mathcal G^{\ell'\ell1}_{m'mm''}Y_{\ell'm'}$.
In the frame $\Khat=\hat{\mathbf z}$, $\vk_{1,2}$ in the $x$--$z$ plane ($\cos\alpha_1=(k_1+k_2\mu)/K$, $\cos\alpha_2=(k_2+k_1\mu)/K$,
$\mu=\khat_1\cdot\khat_2$), $K_{2,\ell m}\to\tilde K_{\ell m}(k_1,k_2,\mu)=\int\dd\eta\,K\sum_L i^L\sqrt{4\pi(2L+1)}j_L(K\chi)\langle\ell m|\mathcal Q|L0\rangle$
with $|m|\le d$, and isotropy reduces the angular integrals of \eqref{eq:loops} to
\begin{align}
\text{(C1)}&\quad C^{(22)}_\ell=\frac1{4\pi^4(2\ell+1)}\int k_1^2\dd k_1\,k_2^2\dd k_2\,P(k_1)P(k_2)\int_{-1}^1\!\dd\mu\sum_m\big|\tilde K_{\ell m}(k_1,k_2,\mu)\big|^2,\nonumber\\
\text{(C2)}&\quad \frac1{2\ell+1}\sum_m\langle C_{\ell m}B^*_{\ell m}\rangle=\frac1{4\pi^4(2\ell+1)}\int k_1^2\dd k_1k_2^2\dd k_2\,PP\int\!\dd\mu\sum_m\tilde K^C_{\ell m}\tilde K^{B*}_{\ell m},
\label{eq:C1C2}
\end{align}
where for a type-B map $\tilde K^B_{\ell m}=\tfrac12\sum\mathcal G^{\ell\ell_1\ell_2}_{mm_1m_2}[\tilde A_{\ell_1m_1}(k_1)\tilde A'_{\ell_2m_2}(k_2)+(A\leftrightarrow A')]$,
$\tilde A_{\ell_1m_1}(k)=4\pi\,i^{\ell_1}Y^*_{\ell_1m_1}(\khat)\Delta^A_{\ell_1}(k)$ at the frame angles.  A non-symmetric
kernel must be symmetrised as $\tfrac12(K+K^{\rm T})$ \emph{before} the $m$-sum.  Two-time kernels use the
double Rayleigh expansion instead of \eqref{eq:K2lm}; (C1) is unchanged.

\subsection{The \texorpdfstring{$(13)$}{(13)} term}
\label{sec:thirteen}
From \eqref{eq:theta_orders},
\begin{equation}
\text{(E0)}\qquad C^{(13)}_\ell=T_{3,\ell}-3T_{\sky1 2,\ell}-3T_{\sky2 1,\ell}+6\langle\sky{m_1}^2\rangle C^{(11)}_\ell,\qquad
T_{3,\ell}\equiv\tfrac1{2\ell+1}\textstyle\sum_m\langle m_{1,\ell m}m^*_{3,\ell m}\rangle,
\label{eq:E0}
\end{equation}
and $T_{\sky12}$, $T_{\sky21}$ the corresponding cross terms with $\sky{m_1}m_2$, $\sky{m_2}m_1$.

\paragraph{The contracted map and $\Delta^{(3c)}_\ell$ (defined once).}  The $q$-contracted third-order map is
\begin{equation}
\Theta^{(3)}_c(\vk;\nhat)\equiv3\int\!\frac{\dd^3q}{(2\pi)^3}P(q)\,K_3(\vk,\vq,-\vq;\nhat),\qquad
C^{(13)}_\ell=\frac1{2\ell+1}\sum_m\int\!\frac{\dd^3k}{(2\pi)^3}P(k)\,K_{1,\ell m}(\vk)\,\Theta^{(3)*}_{c,\ell m}(\vk),
\label{eq:Theta3c}
\end{equation}
i.e.\ the pairing factor 3 of \eqref{eq:loops} is included in $\Theta^{(3)}_c$ once and for all.  After the $\qhat$ integration it depends on
$\nhat$ only through $\nu\equiv\nhat\cdot\khat$; its part with all fields at one point of the ray (structure (iii) below) is written as
\begin{equation}
\Theta^{(3)}_c(\vk;\nhat)\supset\int\!\dd\eta\sum_Lh_L(k,\eta)\,P_L(\nhat\cdot\khat)\,e^{i\vk\cdot\nhat\chi},\qquad
h_L(k,\eta)=\int\!\frac{q^2\dd q}{2\pi^2}P(q)\,h_L(k,q,\eta),
\label{eq:hL}
\end{equation}
with $h_L(k,q,\eta)$ the Legendre coefficients in $\nu$ of the $\qhat$-averaged kernel, $h_L=\frac{2L+1}2\int_{-1}^1\dd\nu\,P_L(\nu)\,F(\nu)$ if
the $\qhat$-average of the integrand is $F(\nu)e^{ik\nu\chi}$.  Every such Legendre series is handled by
\begin{equation}
\text{(A-$L$)}\qquad h(k,\eta)P_L(\nhat\cdot\khat)e^{i\vk\cdot\nhat\chi}\ \Rightarrow\
\Delta_\ell=\int\!\dd\eta\,h\sum_{J=|\ell-L|}^{\ell+L}(2J+1)\threej{\ell}{L}{J}{0}{0}{0}^2i^{J-\ell}j_J(k\chi),
\label{eq:AL}
\end{equation}
and the transfer function of the contracted map is
\begin{equation}
\Delta^{(3c)}_\ell(k)\equiv\int\!\dd\eta\sum_Lh_L(k,\eta)\sum_{J}(2J+1)\threej{\ell}{L}{J}{0}{0}{0}^2i^{J-\ell}j_J(k\chi),
\qquad \Theta^{(3)}_{c,\ell m}=4\pi\,i^\ell Y^*_{\ell m}(\khat)\,\Delta^{(3c)}_\ell(k).
\label{eq:Delta3c}
\end{equation}
The $\vq$-angular reduction of products of fields at two points of the ray is
\begin{equation}
\text{(E1)}\qquad \int\!\frac{\dd\Omega_q}{4\pi}g(\qhat\cdot\khat)(\nhat\cdot\qhat)^be^{i\vq\cdot\nhat\Delta\chi}
=\sum_{L'}i^{L'}(2L'+1)j_{L'}(q\Delta\chi)\sum_{L''}\beta_{L'bL''}\,g_{L''}\,P_{L''}(\nhat\cdot\khat),
\label{eq:E1}
\end{equation}
with $g_{L''}=\tfrac12\int\dd\mu\,P_{L''}(\mu)g(\mu)$ and $x^bP_{L'}=\sum_{L''}\beta_{L'bL''}P_{L''}$; for $\Delta\chi=0$ only
$L'=0$ survives.  The three structures in $m_3$ are then:
\begin{itemize}
\item[(iii)] genuine third-order kernels (local in the ray): $T^{\rm(iii)}_{3,\ell}=\frac2\pi\int k^2\dd k\,P\,\Delta_\ell(k)\Delta^{(3c)}_\ell(k)^*$;
\item[(i)] products of three type-A maps $ABC$: $T^{\rm(i)}_{3,\ell}=C^{m_1A}_\ell\xi^{BC}(0)+C^{m_1B}_\ell\xi^{AC}(0)+C^{m_1C}_\ell\xi^{AB}(0)$, with
$\xi^{XY}(0)$ in the closed form \eqref{eq:xi0} (and $\xi^{\nabla B\cdot\nabla C}(0)$ likewise; $\nabla B\nabla C\!:\!\nabla\nabla A\to-\ell(\ell+1)R_{BC}C^{m_1A}_\ell$,
$R_{BC}=\tfrac12\xi^{\nabla B\cdot\nabla C}(0)$);
\item[(ii)] type-A $\times$ type-C products $A\cdot C$: $T^{\rm(ii)}_{3,\ell}=C^{m_1A}_\ell\langle\!\langle C\rangle\!\rangle+2\cdot\frac2\pi\int k^2\dd k\,P\,\Delta^{m_1}_\ell\Delta^{AC*}_\ell$,
$\langle\!\langle C\rangle\!\rangle=\int\frac{\dd^3k}{(2\pi)^3}P(k)K^C(\vk,-\vk;\nhat)$ and $\Delta^{AC}_\ell$ the transfer function of
$\int\frac{\dd^3q}{(2\pi)^3}P(q)K^A(\vq;\nhat)K^C(\vk,-\vq;\nhat)$ obtained from (E1) with $\Delta\chi=\chi_A-\chi_C$ and (A-$L$).
\end{itemize}
The normalisation pieces, with $\bar K_2(\vq_1,\vq_2)\equiv K_{2,00}(\vq_1,\vq_2)/\sqrt{4\pi}$ the sky-average kernel of $m_2$, are
\begin{align}
\text{(E2)}&\quad T_{\sky21,\ell}=\langle\!\langle m_2\rangle\!\rangle C^{(11)}_\ell+\frac1{\pi^3}\int k^2\dd k\,{k'}^2\dd k'\,PP\,\Delta_\ell(k)\Delta_\ell(k')^*\int\!\dd\mu\,P_\ell(\mu)\bar K_2(k\khat,-k'\khat')^*,
\quad \langle\!\langle m_2\rangle\!\rangle=\frac1{2\pi^2}\int k^2\dd k\,P\,\bar K_2(\vk,-\vk),\nonumber\\
\text{(E3)}&\quad T_{\sky12,\ell}=\frac2{2\pi^2}\int q^2\dd q\,P(q)\Delta_0(q)\,\frac2\pi\int k^2\dd k\,P(k)\Delta_\ell(k)\bar\Delta^{m_2}_\ell(k;q)^*,\qquad
\bar\Delta_\ell(k;q)=\frac{\tfrac12\int\dd\mu\,K_{2,\ell0}(k\hat{\mathbf z},\vq)}{i^\ell\sqrt{4\pi(2\ell+1)}},
\label{eq:E2E3}
\end{align}
and the same objects give the cross term of \eqref{eq:C22_B}.  The monochromatic $C^{(13)}_\ell$ is computed directly from $\kappa$
(Sec.~\ref{sec:asm-K3c}) and does not need these formulas, but (E2) and (E3) were tested against the corresponding normalisation channels of $\kappa$:
for $\ell=1,\dots,6$ they reproduce the channels $-3\sky{m_2}m_1$ ($+3.96321$ in $\Delta^\kappa_2$) and $-3\sky{m_1}m_2$ ($-0.59274$) to $10^{-6}$ and $2\times10^{-5}$
respectively, the accuracy of the $\mu$ quadrature (H-14).  $T^{\rm(ii)}$ is derived by Wick counting in the same way and has not been tested
numerically; (E3) also collapses to \eqref{eq:monopole_terms}(ii) for a type-B $m_2$ (H-10).  For our map the decomposition into $h_L$, $K^C$ is not
displayed: $\kappa=\Theta^{(3)}_c$ is assembled directly from the term lists by polarisation (Sec.~\ref{sec:asm-K3c}), with its class decomposition
\eqref{eq:asm-C13split}.

\subsection{Control case: lensed \texorpdfstring{$C_\ell$}{Cl}}
\label{sec:lensing_control}
For the remapping $\tilde\Theta=\Theta+\nabla\phi\cdot\nabla\Theta+\tfrac12\nabla^i\phi\nabla^j\phi\nabla_i\nabla_j\Theta$ with
$\Theta=\eps\Theta_1$, $\phi=\eps\phi_1$ one has $\Theta^{(2)}=2\nabla\phi\cdot\nabla\Theta$, $\Theta^{(3)}=3\nabla\phi\nabla\phi\!:\!\nabla\nabla\Theta$,
and the bookkeeping \eqref{eq:cl_target} with (B1) (gradient weight \eqref{eq:wgrad}) and the (i)-type pairing gives exactly
\begin{equation}
\tfrac14C^{(22)}_\ell=\sum_{\ell_1\ell_2}\frac{(2\ell_1+1)(2\ell_2+1)}{4\pi}\threej{\ell}{\ell_1}{\ell_2}{0}{0}{0}^2w_\nabla(\ell,\ell_1,\ell_2)^2\Big[C^{\phi\phi}_{\ell_1}C^{\Theta\Theta}_{\ell_2}+C^{\phi\Theta}_{\ell_1}C^{\Theta\phi}_{\ell_2}\Big],\qquad
\tfrac13\Re C^{(13)}_\ell=-\ell(\ell+1)R\,C^{\Theta\Theta}_\ell,
\end{equation}
$R=\tfrac12\langle|\nabla\phi|^2\rangle$.  Because $w_\nabla$ is symmetric in $\ell_1\leftrightarrow\ell_2$, the correlated cross term
$C^{\phi\Theta}C^{\Theta\phi}$ carries the same weight as the uncorrelated one.  For uncorrelated $\phi,\Theta$ these are Hu's (62)--(63)
\cite{Hu2000,LewisChallinor}, with $F_{\ell\ell_1\ell_2}=w_\nabla\sqrt{(2\ell+1)(2\ell_1+1)(2\ell_2+1)/4\pi}\,(\begin{smallmatrix}\ell&\ell_1&\ell_2\\0&0&0\end{smallmatrix})$;
the cross term is the extension to correlated $\phi,\Theta$, relevant here because the deflection correlates with SW/ISW.  This was
verified algebraically and by Monte Carlo on the sphere (App.~\ref{app:ledger}, H-5), validating the prefactors of \eqref{eq:loops},
the two-pairing structure of (B1) and the monopole$\times$map rules.

\subsection{Monochromatic reduction: definitions}
\label{sec:monodef}
The monochromatic primordial spectrum is defined with the variance $\langle\zeta_1^2\rangle=\sigma^2$ held fixed as $k_0$ varies:
\begin{equation}
P(k)=\frac{2\pi^2\sigma^2}{k_0^2}\,\delta(k-k_0)\quad\Longleftrightarrow\quad\langle\zeta_1^2(\vx)\rangle=\sigma^2,\qquad
\int k^2\dd k\,P\,F(k)\to2\pi^2\sigma^2F(k_0),\quad\int\!\frac{\dd^3k}{(2\pi)^3}P\,F(\vk)\to\sigma^2\langle F(k_0\khat)\rangle_{\khat}.
\label{eq:mono}
\end{equation}
It is applied only after the full observable has been assembled for general $P(k)$.  The expansion parameter is then $\eps^2\sigma^2$.
Under \eqref{eq:mono} the general formulas collapse to (all transfer functions at $k_0$):
\begin{align}
&C^{AA'}_\ell=4\pi\sigma^2\Delta^A_\ell\Delta^{A'*}_\ell,\qquad \langle\sky A\sky{A'}\rangle=\sigma^2\Delta^A_0\Delta^{A'}_0,\qquad
\xi^{XY}(0)=\sigma^2\,\tfrac12\int_{-1}^1\dd\nu\,x(k_0,\nu)y(k_0,\nu)^*;\nonumber\\
&\text{(B1)}:\ C^{BB'}_\ell=(4\pi\sigma^2)^2\sum_{\ell_1\ell_2}\frac{(2\ell_1+1)(2\ell_2+1)}{4\pi}\threej{\ell}{\ell_1}{\ell_2}{0}{0}{0}^2ww'
\big[\Delta^A_{\ell_1}\Delta^C_{\ell_1}\Delta^{A'}_{\ell_2}\Delta^D_{\ell_2}+\Delta^A_{\ell_1}\Delta^D_{\ell_1}\Delta^{A'}_{\ell_2}\Delta^C_{\ell_2}\big];\nonumber\\
&\text{(C1)}:\ C^{(22)}_\ell=\frac{\sigma^4}{2\ell+1}\int_{-1}^1\!\dd\mu\sum_{|m|\le d}\big|\tilde K_{\ell m}(k_0,k_0,\mu)\big|^2,\quad K=k_0\sqrt{2(1+\mu)},\ \cos\alpha_1=\cos\alpha_2=\sqrt{(1+\mu)/2};\nonumber\\
&\text{(13)(iii)}:\ T^{\rm(iii)}_{3,\ell}=4\pi\sigma^2\,\Delta_\ell(k_0)\,\Delta^{(3c)}_\ell(k_0)^*,\qquad
\Delta^{(3c)}_\ell(k_0)=\sigma^2\int\!\dd\eta\sum_Lh_L(k_0,k_0,\eta)\sum_J(2J+1)\threej{\ell}{L}{J}{0}{0}{0}^2i^{J-\ell}j_J(k_0\chi),\nonumber\\
&\langle\sky{m_1}^2\rangle=\sigma^2\Delta_0^2,\quad \langle\!\langle m_2\rangle\!\rangle=\sigma^2\bar K_2(k_0\khat,-k_0\khat),\quad
\text{(E2) cross}=4\pi\sigma^4|\Delta_\ell|^2\!\int\!\dd\mu\,P_\ell(\mu)\bar K_2(k_0\khat,-k_0\khat')^*,\quad \text{(E3)}=8\pi\sigma^4\Delta_0\Delta_\ell\bar\Delta_\ell(k_0;k_0)^* ,
\label{eq:mono_rules}
\end{align}
where in (13)(iii) the pairing factor 3 is inside $\Delta^{(3c)}_\ell$ by \eqref{eq:Theta3c} and $|\vk\pm\vq|=k_0\sqrt{2(1\pm\mu)}$.

\paragraph{Convergence of the $\ell_1,\ell_2$ sums (analytic bound).}  For the monochromatic spectrum every transfer function is a finite
combination of $j^{(s)}_\ell(k_0\chi)$ with $0\le\chi\le\chi_{\max}=\chi_*$ (endpoint and line-of-sight terms; the observer terms only touch
$\ell\le1$).  From $j_\ell(x)=\frac{x^\ell}{2^{\ell+1}\ell!}\int_{-1}^1e^{ixt}(1-t^2)^\ell\dd t$ one has, for real $x\ge0$,
\begin{equation}
|j_\ell(x)|\le\frac{x^\ell}{(2\ell+1)!!}\le\Big(\frac{e\,x}{2\ell+1}\Big)^\ell,\qquad
|j_\ell'(x)|\le\Big(\frac{\ell}{x}+1\Big)\frac{x^\ell}{(2\ell+1)!!},
\label{eq:jbound}
\end{equation}
using $\ln(2\ell+1)!!\ge\int_0^\ell\ln(2t+1)\dd t\ge\ell\ln[(2\ell+1)/e]$.  Hence, with $X\equiv k_0\chi_*$ and $\mathcal N_A$ the total absolute weight of the
map ($\sum_s\int|g_s|\dd\eta$ with the factor $(\ell/X+1)$ for $s=1$), $|\Delta^A_\ell(k_0)|\le\mathcal N_A(\ell/X+1)(eX/(2\ell+1))^\ell$, which for
$\ell\ge\ell_c\equiv\lceil eX\rceil$ is $\le\mathcal N_A(\ell/X+1)\,2^{-\ell}$: the decay is super-exponential.  In (B1) the sum over $\ell_2$ at
fixed $\ell,\ell_1$ is bounded with $\sum_{\ell_2}(2\ell_2+1)(\begin{smallmatrix}\ell&\ell_1&\ell_2\\0&0&0\end{smallmatrix})^2=1$, and $|w|\le(\ell_1+\ell+1)^2$
for $|\ell_1-\ell_2|\le\ell$, while $|\Delta_\ell|\le\mathcal N$ for all $\ell$ ($|j_\ell|,|j'_\ell|\le1$).  Denoting by $C^{BB'}_\ell[L]$ the sum truncated at
$\ell_1,\ell_2\le L$, a geometric-series estimate of the omitted terms gives, for $L\ge\max(\ell_c,10)+\ell$,
\begin{equation}
\big|C^{BB'}_\ell-C^{BB'}_\ell[L]\big|\le8\pi\sigma^4\,\mathcal N_A\mathcal N_{A'}\mathcal N_C\mathcal N_D\,(L+\ell+1)^{6}\,(L/X+1)^2\,4^{\ell-L},
\label{eq:tailbound}
\end{equation}
so the sums are finite to any prescribed accuracy with an explicit, a-priori cut-off; no numerical convergence criterion is used.  The
same bound applies to the type-D shell maps, and the type-C sums \eqref{eq:K2lm} are finite ($|L-\ell|\le d$).

\paragraph{Delicate points} (treated in Secs.~\ref{sec:assembly} and \ref{sec:mono}).
(1)~$\mu\to-1$ in (C1), where $K\to0$, $\dd\mu=K\dd K/k_0^2$ and $j_L(K\chi)\to\delta_{L0}$: a kernel $\sim K^{-p}$ converges only for $p<1$, so
every $1/K^2$ (Poisson inversion) must come with a vanishing numerator (treated in Sec.~\ref{sec:mono-K0}: bounded for $\ell\ge2$).
(2)~$\mu\to\pm1$ in the $(13)$ loop, where $|\vk\mp\vq|\to0$ inside $K_3(\vk,\vq,-\vq)$ and $C_3$; at $q=k$, $C_3$ has a simple pole
\eqref{eq:C3pole}, and the monochromatic $(13)$ integral exists because the pole cancels in $\Theta^{(3)}_c$ (Proposition~\ref{prop:mono-residues}).
(3)~The $K=0$ tadpoles $\langle\!\langle m_2\rangle\!\rangle$, $\langle\!\langle C\rangle\!\rangle$ require kernels exactly at $\vK=0$, where the $k\to0$ limit does
not exist without angular averaging \eqref{eq:KPhi20_smallk}; they enter only through \eqref{eq:theta_orders}, where the sky-mean
normalisation cancels the direction-independent mean of $e^\Lambda$, and are kept symbolic until combined; the $\qhat$-averaged limit exists and is the correct one (Sec.~\ref{sec:asm-zeromodes}).
(4)~$\langle\zeta_{\vk_1}\zeta_{\vk_2}\rangle$ inside $\Theta^{(2)}$ feeds only $\ell=0$, consistent with the monopole normalisation.
